\documentclass[11pt]{article}
\usepackage[letterpaper,landscape,margin=0.65in]{geometry}
\usepackage{amsmath}

% Editable spacing. Matrix cells are square; vectors use one column.
\newcommand{\CellPadding}{4pt}
\newcommand{\CellHeightFactor}{1.4}
\newcommand{\DividerPadding}{2pt}

% Working sizes are local to each box.
\newsavebox{\CellSample}
\newlength{\CellSide}
\newlength{\WritingWidth}
\newlength{\WritingHeight}
\newlength{\LeftBlockWidth}

% Center on the math axis so parentheses fit the actual writing height.
\newcommand{\WritingRegion}[2]{%
  \vcenter{\hbox{\rule{0pt}{#2}\makebox[#1]{}}}}

% Same arguments: [style]{rows}{columns}{points}{sample text}.
% Measure one invisible sample. Square cells replace phantom arrays.
\newcommand{\AnswerBox}[5][plain]{%
  \begingroup
  \sbox{\CellSample}{\fontsize{#4}{#4}\selectfont #5}%
  \setlength{\CellSide}{#4pt}%
  \setlength{\CellSide}{\CellHeightFactor\CellSide}%
  \ifdim\wd\CellSample>\CellSide
    \setlength{\CellSide}{\wd\CellSample}%
  \fi
  \addtolength{\CellSide}{\CellPadding}%
  \setlength{\WritingWidth}{#3\CellSide}%
  \setlength{\WritingHeight}{#2\CellSide}%
  \def\BoxStyle{#1}%
  \def\PlainStyle{plain}%
  \def\AugmentedStyle{augmented}%
  \ensuremath{\boxed{%
    \ifx\BoxStyle\PlainStyle
      \WritingRegion{\WritingWidth}{\WritingHeight}%
    \else
      \left(
      \ifx\BoxStyle\AugmentedStyle
        \setlength{\LeftBlockWidth}{\WritingWidth}%
        \addtolength{\LeftBlockWidth}{-\CellSide}%
        \setlength{\arraycolsep}{\DividerPadding}%
        \begin{array}{@{}c|c@{}}
          \WritingRegion{\LeftBlockWidth}{\WritingHeight}&
          \WritingRegion{\CellSide}{\WritingHeight}
        \end{array}%
      \else
        \WritingRegion{\WritingWidth}{\WritingHeight}%
      \fi
      \right)
    \fi
  }}%
  \endgroup}

\newcommand{\MatrixAnswerBox}[4]{\AnswerBox[matrix]{#1}{#2}{#3}{#4}}
% Columns include the RHS. A square n-by-n left block needs n+1 total columns.
\newcommand{\AugmentedAnswerBox}[4]{\AnswerBox[augmented]{#1}{#2}{#3}{#4}}
\newcommand{\VectorAnswerBox}[3]{\MatrixAnswerBox{#1}{1}{#2}{#3}}

\pagestyle{empty}
\begin{document}
\section*{Worksheet answer boxes}
Use \texttt{\{rows\}\{columns\}\{points\}\{sample text\}}.
Cell size is the larger of the sample width and the scaled point size, plus padding.
The sample stays invisible. Adjust the three spacing settings above to change the look.

\bigskip
\begin{tabular}{@{}lccc@{}}
  & Plain & Parentheses & Augmented\\[6pt]
  1 row & \AnswerBox{1}{2}{24}{00}
        & \MatrixAnswerBox{1}{2}{24}{00}
        & \AugmentedAnswerBox{1}{2}{24}{00}\\[12pt]
  2 rows & \AnswerBox{2}{3}{24}{00}
         & \MatrixAnswerBox{2}{3}{24}{00}
         & \AugmentedAnswerBox{2}{3}{24}{00}\\[12pt]
  3 rows & \AnswerBox{3}{3}{24}{00}
         & \MatrixAnswerBox{3}{3}{24}{00}
         & \AugmentedAnswerBox{3}{4}{24}{00}
\end{tabular}

\bigskip
\verb|\MatrixAnswerBox{3}{3}{24}{00}| gives a square matrix writing area.\par
\verb|\VectorAnswerBox{3}{24}{00}| has the same height and one-column width.\par
\verb|\AugmentedAnswerBox{3}{4}{24}{00}| adds one RHS column to a square left block.

\medskip
Boxes use standard \verb|\boxed| padding.
Parentheses surround the writing region; the augmented divider stays visible.
\end{document}
